\section{离散状态的扩散}
前面的空间分数使用 $\nabla_x$，要求连续坐标上的可微密度。
但这不意味着离散数据一定要先经 VAE 才能做扩散。
对于词表、类别等有限状态，可以直接规定随机替换或遮蔽的马尔可夫过程，
再学习相应的逆转移。这里用概率比值代替空间梯度。

先看一个小时间步。设从状态 $i$ 跳到 $j$ 的速率为 $Q_t(i,j)$，$i\ne j$，
并令 $Q_t(i,i)=-\sum_{j\ne i}Q_t(i,j)$。
若 $p_t$ 写为行向量，主方程就是 $\partial_tp_t=p_tQ_t$。
对逆向时钟，贝叶斯公式给出从当前 $j$ 跳回 $i$ 的速率
\[
\bar Q_t(j,i)=Q_t(i,j)\frac{p_t(i)}{p_t(j)},\qquad p_t(j)>0.
\]
推导只需将
$P(X_{t-h}=i\mid X_t=j)$ 写为正向转移概率乘以前一时刻边缘、再除以当前边缘，
对 $i\ne j$ 除以 $h$ 并取极限。
SEDD 学习的离散分数正是这样的密度比值，而不是对 token 编号求导。\cite{sedd}

掩码扩散是更具体的例子。假设一个 token 到时间 $t$ 仍保留的概率为 $\rho_t$，
否则变成吸收态 $\mathtt{MASK}$。对 $s<t$ 且 $\rho_s>\rho_t$，
在已知 $X_t=\mathtt{MASK}$ 时，逆向重新显露原 token 的概率为
\[
P(X_s\ne\mathtt{MASK}\mid X_t=\mathtt{MASK})
=\frac{\rho_s-\rho_t}{1-\rho_t}.
\]
分子是“在 $s$ 尚未遮蔽、到 $t$ 已被遮蔽”的概率，分母是当前已遮蔽概率。
至于显露成哪个 token，则由模型在上下文与时间条件下预测，
相应训练常写成带时间权重的遮蔽位置交叉熵。
多位置迭代可以并行更新，但同一步的条件独立预测、有限步采样和具体重遮蔽策略仍会影响结果。
它与自回归逐 token 分解、连续高斯去噪有联系，也各有自己的逆过程定义。\cite{mdlm,theorybook}
